\section{连续缓存分配与换入换出}
\subsection{驻留段约束}
缓存容量为 $C_{L1}=4096$、$C_{UB}=1024$、$C_{L0A}=C_{L0B}=256$、$C_{L0C}=512$。对缓冲区 $b$ 的每个驻留段 $j$，要求
\begin{equation}0\le x_{bj},\qquad x_{bj}+s_b\le C_{t_b}.\end{equation}
若两个同类型驻留段在调度序列中相交，则必须满足
\begin{equation}
x_{bj}+s_b\le x_{ak}\quad\text{或}\quad x_{ak}+s_a\le x_{bj}.
\end{equation}
后一条件是连续地址分配的关键，容量总和约束本身不能排除碎片。

记 $\chi_b=1$ 表示缓冲区被输入复制节点使用，否则为零。一次换出的额外搬运量按题面为 $(2-\chi_b)s_b$。因此目标为
\begin{equation}D=\sum_{(b,j)\in\mathcal{M}}(2-\chi_b)s_b,\end{equation}
其中 $\mathcal{M}$ 为所有换出事件的集合，重复换出必须重复计费。换入时间统一为 $2s_b+150$；换出时间在 $\chi_b=1$ 时为零，否则同为 $2s_b+150$。

\subsection{分配和牺牲对象选择}
沿问题一候选序列维护各类缓存空闲区间。新申请先选择能容纳它的最小空闲区间，以降低剩余碎片。空间不足时，排除当前操作引用的缓冲区，从同类型驻留对象中选择牺牲者。第一种准则最大化距离下一次使用的序列距离 $\Delta_b$；第二种最大化 $\Delta_b/[(2-\chi_b)s_b]$，以兼顾未来使用距离与搬运代价。无后续使用但尚不能逻辑释放的对象取序列尾部之后的距离，仍保留其配对换入与最终释放。

换出后立即释放该驻留段的地址，在后续使用前重新申请连续区间并插入换入节点。换入允许采用不同偏移。每次换出对应一对新增节点，原始节点总数为 $N$ 时，第 $j$ 次的编号为 $N+2j$、$N+2j+1$，其中 $j$ 从零开始。

若受保护缓冲区本身分散在多个地址区间，使其他所需对象无法放入，即使总量不超容量，单纯换出非保护对象也可能失败。此时将当前所需缓存类型的驻留对象显式换出，再按大小递减重装本次操作的全部对象。清空后的连续空间保证容量检查通过的操作能够重新装入；代价可能较大，故仅作为可行性恢复机制。禁止将此过程处理成无成本的内存整理。

算法每次空间申请扫描缓存单元并比较牺牲者，开销受缓存长度和换出次数共同控制。以 $C=\max_t C_t$、$M$ 为换出次数，其分配和牺牲者扫描可保守界为 $O((N+M)(C+|B|))$；未来使用位置通过队列预先索引。题目给定容量固定，但一般化分析中不能忽略 $C$。

\subsection{结果与物理验证}
问题二允许调整问题一的顺序，因此在编号、缓存贪心和深度优先的最佳适配候选中，先选择搬运量最小者，再以时间打破平局。其结果不必与问题一最小峰值候选相同。
\begin{table}[H]\centering\caption{缓存优先方案的额外搬运量与执行时间}\small
\begin{tabular}{lrrr}\toprule 计算图 & 搬运量 & 换出次数 & 周期\\\midrule
Conv\_Case0 & 44766 & 56 & 472038 \\
Conv\_Case1 & 75546 & 2116 & 1210619 \\
FlashAttention\_Case0 & 3620 & 23 & 47509 \\
FlashAttention\_Case1 & 32852 & 212 & 225902 \\
Matmul\_Case0 & 34944 & 273 & 194331 \\
Matmul\_Case1 & 460800 & 3600 & 1800001 \\
\bottomrule\end{tabular}
\end{table}
所有方案通过逐单元的地址占用检查，确保没有越界、重叠或非驻留使用。矩阵乘大图的搬运量仍然较大，说明有限缓存下的大量复用数据不能长期驻留。当前方法没有给出最小搬运量下界，因此不能将此数值解释为无法进一步降低的必要开销。

